% problem-000300 1 硅与超强碳正离子
\section*{1. 硅与超强碳正离子}
\textit{下列为分问。}

\subsection*{1-1}
已知反应分两步进⾏，试⽤化学⽅程式表⽰上述溶解过程。

早在上世纪50年代就发现了 $\mathrm { C H } _ { 5 } { } ^ { + }$ 的存在，⼈们曾提出该离⼦结构的多种假设，然⽽，直⾄1999年，才在低温下获得该离⼦的振动-转动光谱，并由此提出该离⼦的如下结构模型：氢原⼦围绕着碳原⼦快速转动；所有C—H键的键⻓相等。

\subsection*{1-2}
该离⼦的结构能否⽤经典的共价键理论说明？简述理由。

\subsection*{1-3}
该离⼦是（ ）。A 质⼦酸 B 路易斯酸 C ⾃由基 D 亲核试剂

2003年5⽉报道，在⽯油中发现了⼀种新的烷烃分⼦，因其结构类似于⾦刚⽯，被称为“分⼦钻⽯”，若能合成，有可能⽤做合成纳⽶材料的理想模板。该分⼦的结构简图如下：

（图：）

\subsection*{1-4}
该分⼦的分⼦式为 ？1-5 该分⼦有⽆对称中⼼？1-6 该分⼦有⼏种不同级的碳原⼦？1-7 该分⼦有⽆⼿性碳原⼦？1-8 该分⼦有⽆⼿性？
